\documentclass[11pt]{article}
\usepackage[margin=1in]{geometry}
\usepackage{amsmath,amssymb,amsthm,booktabs,array}
\usepackage[T1]{fontenc}
\usepackage{xcolor}
\usepackage[colorlinks=true,linkcolor=blue!60!black,urlcolor=blue!60!black]{hyperref}
\newcommand{\E}{\mathbb E}
\newcommand{\R}{\mathbb R}
\newcommand{\ip}[2]{\langle #1,#2\rangle}
\newcommand{\norm}[1]{\lVert #1\rVert}
\newcommand{\TV}{\mathrm{TV}}
\newcommand{\on}{\mathrm{on}}
\newcommand{\exit}{\mathrm{exit}}
\newcommand{\loc}{\mathrm{loc}}
\newcommand{\He}{\mathrm{He}}
\newtheorem{theorem}{Theorem}[section]
\newtheorem{proposition}[theorem]{Proposition}
\newtheorem{lemma}[theorem]{Lemma}
\theoremstyle{remark}
\newtheorem{remark}[theorem]{Remark}
\title{Correlated Features Beyond Hermite Polynomials\\
\large An exact Gegenbauer calculation, a Legendre exception,\\and implications for neural-network students}
\author{Research note for the Correlated Features project}
\date{October 8, 2026}
\begin{document}
\maketitle

\begin{abstract}
We redo the feature-overlap, recruitment, and shared-population stability
calculations with Gegenbauer ridge functions on uniformly distributed
spherical inputs. At the level of the fixed-penalty population objective,
the two-teacher Hermite picture survives in every ambient dimension
$D\geq4$, with explicit modified thresholds and an all-direction
optimality certificate. In $D=3$, the functions are Legendre polynomials:
an additional regime of directions outside the teacher plane intervenes,
and the shared state loses optimality before its in-plane splitting
instability. On the circle, a single Fourier frequency needs no splitting
at all. We also solve a Gaussian linear--cubic mixture and give an exact
obstruction to learning cubic teachers with bias-free ReLU students.
The overlap and optimization results below are exact. Transferring the
microscopic Bayesian posterior, its background integration, and its neuron
counts to these new models remains a separate problem.
\end{abstract}

\section{What is being changed, and what is being proved}

The source paper's Appendix E and the correlated-features notes reduce the
cubic problem to
\begin{equation}\label{eq:original}
 F=\lambda\sum_j b_j+\frac12\E(f_*-f_C)^2,
 \qquad f_C=\sum_j b_j\psi_3(u_j^\top X),\quad b_j\geq0,
 \qquad \psi_3(z)=\frac{z^3-3z}{\sqrt6}.
\end{equation}
Here $X$ is Gaussian, unit directions have overlap kernel $q^3$, and
$\lambda$ incorporates the alignment cost and the fixed effective data
force. The present note asks which conclusions follow from this kernel,
and which follow from the more general fit-versus-population-cost structure.

We hold the form of the population penalty fixed while changing the data
measure and the normalized feature family. This gives an exact comparison
of the reduced variational models. It does \emph{not} assert that the original
Bayesian network induces the same penalty at finite $D$. Section~\ref{sec:micro}
rederives the relevant geometric quantities and explains that limitation.
All comparisons of signal use $\beta/\lambda$ within the specified model.

Two dimensions must be distinguished. $D$ is the ambient input dimension;
$m$ is the dimension of the teacher span. The earlier harmonic decomposition
of the Hermite overlap on $S^{m-1}$ did not change the input law. Here the
input law itself changes, and the new overlap kernel depends on $D$, even
when there are only two teachers and $m=2$.

\paragraph{Why orthogonality is not enough.}
For a scalar marginal law $\nu$, orthogonality only gives
$\E_\nu[p_k(Z)p_l(Z)]=\delta_{kl}$. It does not determine
$\E[p_k(u^\top X)p_l(v^\top X)]$ for different directions. The Gaussian
Hermites have the stronger identities
\begin{equation}\label{eq:hermite}
 \E[\psi_k(u^\top X)\psi_l(v^\top X)]=\delta_{kl}(u^\top v)^k,
 \qquad
 \E_Z\psi_k(rz+\sqrt{1-r^2}Z)=r^k\psi_k(z).
\end{equation}
The second identity is what makes a partially aligned group a rescaled copy
of the same ridge function. Both identities need attention when the model
is changed. Legendre polynomials are not orthogonal under a Gaussian law.

\section{The new orthogonal family and its exact overlap}

Let $S$ be uniform on the unit sphere $S^{D-1}\subset\R^D$, with normalized
probability measure, and take $D\geq3$. The marginal $t=u^\top S$ has density
proportional to $(1-t^2)^{(D-3)/2}$ on $[-1,1]$. Define
\begin{equation}
 Z_l^{(D)}(t)=\frac{C_l^{((D-2)/2)}(t)}{C_l^{((D-2)/2)}(1)},
 \quad h_{l,D}=\binom{D+l-1}{l}-\binom{D+l-3}{l-2},
 \quad \phi_{l,u}(S)=\sqrt{h_{l,D}}\,Z_l^{(D)}(u^\top S),
\end{equation}
where the second binomial is zero for $l<2$. These are unit-variance
zonal spherical harmonics for $l\geq1$.

\begin{proposition}[Rotated orthogonality]\label{prop:addition}
For unit $u,v$,
\begin{equation}\label{eq:addition}
 \E[\phi_{l,u}(S)\phi_{k,v}(S)]
   =\delta_{lk}Z_l^{(D)}(u^\top v).
\end{equation}
In particular, for cubic features,
\begin{equation}\label{eq:kernel}
 K_D(q)=Z_3^{(D)}(q)=a_Dq^3-b_Dq,
 \qquad a_D=\frac{D+2}{D-1},\quad b_D=\frac3{D-1},
 \qquad h_{3,D}=\frac{D(D-1)(D+4)}6.
\end{equation}
\end{proposition}
\begin{proof}
Choose a real orthonormal degree-$l$ harmonic basis $Y_{lj}$ for normalized
spherical measure. The addition theorem gives
$\sum_{j=1}^{h_{l,D}}Y_{lj}(u)Y_{lj}(s)=h_{l,D}Z_l^{(D)}(u^\top s)$.
Thus $\phi_{l,u}=h_{l,D}^{-1/2}\sum_jY_{lj}(u)Y_{lj}$.
Taking inner products proves the assertion; different harmonic degrees
are orthogonal. The cubic expression also follows directly from
$\E S_1^2=1/D$, $\E S_1^4=3/[D(D+2)]$, and the sixth moments.
\end{proof}

The addition theorem and Gegenbauer normalization are standard
\cite{dlmf,frye}; the phase calculations below are derived here.
At $D=3$, $K_3=P_3=(5q^3-3q)/2$ and
$\phi_{3,u}=\sqrt7 P_3(u^\top S)$. The \emph{direction overlap} $\omega$
and the \emph{feature correlation} $K_D(\omega)$ are different quantities.
The latter can be negative even when $\omega>0$.

\section{Population cost and a global residual test}

For a normalized feature dictionary $\phi_u$ and an arbitrary square
integrable target, define the student and objective by
\begin{equation}\label{eq:measure}
 f_\mu=\int\phi_u\,d\mu(u),\qquad
 F(\mu)=\lambda\norm{\mu}_{\TV}+\tfrac12\norm{f_*-f_\mu}_{L^2}^2.
\end{equation}
This is also an infinite-width, two-layer ridge network with an $\ell_1$
penalty on output weights; the directions are optimized, not held fixed.
For odd features, $\phi_{-u}=-\phi_u$, so signed weights can be absorbed
into directions and $\norm\mu_\TV$ becomes total nonnegative mass.
This is the variation-norm formulation studied for broader neural-network
families in \cite{bach}.

For $f_*=\beta\sum_{i=1}^k\phi_{e_i}$ with a rotational kernel $K$,
\begin{equation}\label{eq:Fexpand}
 F=\mathrm{constant}+\lambda\sum_j b_j
 -\beta\sum_jb_j H_K(u_j)
 +\tfrac12\sum_{j,l}b_jb_l K(u_j^\top u_l),
 \qquad H_K(u)=\sum_iK(e_i^\top u).
\end{equation}
Thus every teacher--student and student--student interaction must use the
new kernel, rather than changing the target overlap alone.

\begin{proposition}[Global certificate]\label{prop:certificate}
Let $R=f_*-f_\mu$. A nonnegative representation for an odd dictionary
minimizes \eqref{eq:measure} if and only if
\begin{equation}\label{eq:certificate}
 \ip{R}{\phi_u}\leq\lambda\quad\text{for all unit }u,
 \qquad \ip{R}{\phi_u}=\lambda\quad\mu\text{-almost everywhere}.
\end{equation}
Equivalently, $|\ip{R}{\phi_u}|\leq\lambda$ with sign-correct contact
equalities for a signed measure. Recruitment begins at
\begin{equation}\label{eq:onset}
 \beta_\on=\frac\lambda{\max_u H_K(u)}
\end{equation}
when the maximum is positive.
\end{proposition}
\begin{proof}
Adding a small mass $\varepsilon$ at $u$ changes the objective to first
order by $\varepsilon(\lambda-\ip R{\phi_u})$. Removing occupied mass
gives the contact equality. Conversely, for any nonnegative competitor
$\nu$, expanding the difference gives
\[
 F(\nu)-F(\mu)=\tfrac12\norm{f_\nu-f_\mu}^2
 +\int(\lambda-\ip R{\phi_u})\,d\nu(u)
 -\int(\lambda-\ip R{\phi_u})\,d\mu(u)\geq0.
\]
Putting $\mu=0$ proves the onset assertion. No orthogonality assumption
is needed for this argument.
\end{proof}

\section{Two Gegenbauer teachers: reduction over all directions}

Take $e_1^\top e_2=\omega\in[0,1)$ and define
\begin{equation}
 e_{1,2}=ce_+\pm se_-,\qquad
 c=\sqrt{\frac{1+\omega}{2}},\quad
 s=\sqrt{\frac{1-\omega}{2}}.
\end{equation}
Write a student direction as $u=xe_++ye_-+zv$, where
$v\perp\operatorname{span}\{e_1,e_2\}$ and $x^2+y^2+z^2=1$.
All dependence on transverse coordinates is captured by $z^2$.
With $a=a_D,b=b_D$,
\begin{equation}\label{eq:Hxy}
 H_D(x,y)=2cx\{ac^2x^2+3as^2y^2-b\}.
\end{equation}
For $x\geq0$, this is increasing in $y^2$. Its largest value at fixed
$x$ therefore occurs in the teacher plane, $y^2=1-x^2$; its smallest
occurs at $y=0$. Since changing $u$ to $-u$ changes the sign of $H_D$,
\emph{both} extrema must be checked for the global positive maximum.

The planar branch is
\begin{equation}\label{eq:Hplane}
 H_{\mathrm{pl}}(t)=2ct\{A+a(2\omega-1)t^2\},
 \quad A=3as^2-b,\quad 0\leq t\leq1.
\end{equation}
Its positive maximum occurs at
\begin{equation}\label{eq:tstar}
 t_*^2=\frac{3as^2-b}{3a(1-2\omega)}\quad
 \text{if }\omega<\omega_{\mathrm{pl}},\qquad
 t_*=1\quad\text{if }\omega\geq\omega_{\mathrm{pl}},
 \qquad \omega_{\mathrm{pl}}=\frac{D+4}{3(D+2)}.
\end{equation}
The competing positive maximum obtained from $x<0,y=0$ is
\begin{equation}\label{eq:transverse}
 M_\perp=\frac4{(D-1)\sqrt{D+2}},\qquad
 x=-\frac1{c\sqrt{D+2}},\quad y=0.
\end{equation}
This follows by minimizing $2cx(ac^2x^2-b)$ for $x\geq0$.
The minimizing value is feasible for every $\omega\in[0,1)$ and $D\geq3$.

\begin{theorem}[Onset in dimensions at least four]\label{thm:onset}
For $D\geq4$, a global onset maximizer always lies in the teacher plane.
The critical direction correlation is
\begin{equation}\label{eq:wc}
 \boxed{\omega_c(D)=\frac{D+4}{3(D+2)}
       =\frac13+\frac{2}{3(D+2)}.}
\end{equation}
Above it, the maximizing direction is $e_+$. Below it there are the two
reflected directions $t_*e_+\pm\sqrt{1-t_*^2}e_-$, with $t_*$ given by
\eqref{eq:tstar}. Their onset is $\lambda/H_{\mathrm{pl}}(t_*)$.
The two directions merge continuously at $\omega_c(D)$.
\end{theorem}
\begin{proof}
On $[0,1]$, the minimum of $K_D$ is
$-2/[(D-1)\sqrt{D+2}]=-M_\perp/2$. The planar direction $e_1$ therefore
has $H_D(e_1)=1+K_D(\omega)\geq1-M_\perp/2>M_\perp$ for $D\geq4$,
because $(D-1)\sqrt{D+2}>6$. Hence the transverse alternative cannot
win. Differentiating \eqref{eq:Hplane} proves the remaining claims.
\end{proof}

Even orthogonal teacher groups no longer decouple exactly. At
$\omega=0$, the two-teacher onset maximum is
\begin{equation}\label{eq:orthogonal}
 M_D(0)=\frac{D}{D-1}\sqrt{\frac D{D+2}}>1
 \qquad(D\geq3),
\end{equation}
whereas either individual teacher direction has overlap one with the
two-feature target. This follows by substituting $\omega=0$ in
\eqref{eq:tstar} and \eqref{eq:Hplane}; the Legendre case is also verified
by Theorem~\ref{thm:legendre}. Its square exceeds one because
$D^3-(D-1)^2(D+2)=3D-2>0$. Thus $K_D(0)=0$ is insufficient to recover
the orthogonal-block decoupling proved for Hermite cubic atoms. The
discrepancy tends to zero at order $D^{-2}$.

\section{A certified shared interval for every $D\geq4$}

On the shared branch $f_C=m\phi_{e_+}$, normalization gives
\begin{equation}\label{eq:shared}
 m=[2\beta K_D(c)-\lambda]_+,
 \qquad \beta_\on=\frac\lambda{2K_D(c)}.
\end{equation}
For an occupied shared atom, introduce
\begin{equation}\label{eq:B}
 B=\frac{2\beta a c s^2}{\lambda}.
\end{equation}
Its normalized residual correlation, using \eqref{eq:shared}, is exactly
\begin{equation}\label{eq:residual}
 \frac{\ip R{\phi_u}}\lambda
  =(a+bB)x^3+3Bxy^2-b(1+B)x.
\end{equation}
As in the onset calculation, for $x\geq0$ we check the planar maximum
and the transverse minimum. On the plane the residual is
\begin{equation}\label{eq:resplane}
 p_B(x)=\{a-(3-b)B\}x^3+\{(3-b)B-b\}x,\quad p_B(1)=1.
\end{equation}
The local splitting condition is $p_B'(1)=0$, or
\begin{equation}\label{eq:Bcrit}
 B_c=\frac{D+1}{2(D-2)}.
\end{equation}
For $0\leq B\leq B_c$, $p_B(x)\leq1$ on $[0,1]$:
its derivative is affine in $x^2$ and can have at most one interior
minimum of $p_B$ when its initial derivative is negative; if its cubic
coefficient is negative, its derivative is minimized at $x=1$ and is
nonnegative. The endpoint values are $0$ and $1$.

The magnitude of the negative transverse minimum is
\begin{equation}\label{eq:negative}
 C_D(B)=\frac23 b(1+B)
    \sqrt{\frac{b(1+B)}{3(a+bB)}}.
\end{equation}
It is increasing in $B\geq0$, and at the planar threshold
\begin{equation}\label{eq:Ccrit}
 C_D(B_c)=\frac{3\sqrt3}{(D-2)\sqrt{2D+5}}<1
 \quad(D\geq4).
\end{equation}
Thus no out-of-plane atom wins earlier.

\begin{theorem}[Global endpoint of the shared state]\label{thm:exit}
For $D\geq4$ and $\omega>\omega_c(D)$, the shared solution
\eqref{eq:shared} is globally optimal over arbitrary numbers of atoms
and all ambient directions precisely throughout
\begin{equation}\label{eq:bounds}
 \boxed{\frac\lambda{2K_D(c)}\leq\beta\leq
 \beta_\exit=\frac\lambda{4cs^2}\frac{D^2-1}{D^2-4}.}
\end{equation}
At the upper endpoint, an in-plane splitting mode becomes marginal;
for larger $\beta$ the shared solution is not optimal.
\end{theorem}
\begin{proof}
Equations \eqref{eq:resplane}--\eqref{eq:Ccrit} establish the global
certificate for $B\leq B_c$. Above $B_c$, $p_B'(1)<0$, so
$p_B(1-\varepsilon)>1$ for small positive $\varepsilon$, violating it.
Converting $B_c$ to $\beta$ gives \eqref{eq:bounds}. The interval has
positive length exactly when $\omega>\omega_c(D)$; at equality its
endpoints coincide.
\end{proof}

The Hermite limits are $\omega_c\to1/3$, $K_D(c)\to c^3$, and
$\beta_\exit\to\lambda/(4cs^2)$. The correction to the shared-direction
correlation threshold is order $D^{-1}$, whereas the displayed multiplicative
correction to the exit formula at fixed $\omega$ is order $D^{-2}$.
This theorem establishes the onset and the exact shared-branch endpoint.
It does not classify all populations beyond that endpoint.

\section{Legendre polynomials: an additional phase at $D=3$}

At $D=3$, the bound \eqref{eq:Ccrit} exceeds one, so extending the preceding
claim to Legendre polynomials would be incorrect.

\begin{theorem}[Exact Legendre onset]\label{thm:legendre}
For $\phi_u(S)=\sqrt7 P_3(u^\top S)$ and two equal teachers,
the global onset directions are as follows, with coexistence at the boundaries:
\begin{center}
\begin{tabular}{@{}lll@{}}
\toprule
Direction correlation & Maximizing directions & Maximum overlap\\
\midrule
$0\leq\omega<1/5$ & $t_*e_+\pm\sqrt{1-t_*^2}e_-$ & $H_{\mathrm{pl}}(t_*)$\\
$1/5<\omega<3/5$ & $-\dfrac{e_+}{\sqrt5c}\pm
 \sqrt{1-\dfrac1{5c^2}}e_0$ & $2/\sqrt5$\\[5pt]
$3/5<\omega<1$ & $e_+$ & $2P_3(c)$\\
\bottomrule
\end{tabular}
\end{center}
Here $e_0$ is normal to the teacher plane and
\begin{equation}
 t_*^2=\frac{3-5\omega}{10(1-2\omega)}
 \quad(0\leq\omega<1/5).
\end{equation}
Both correlation-driven changes of onset direction are discontinuous.
The middle interval has $\beta_\on=\sqrt5\lambda/2$.
\end{theorem}
\begin{proof}
The all-direction reduction above leaves only the positive planar
maximum and $M_\perp=2/\sqrt5$. The planar stationary branch applies
until $\omega=7/15$. On that branch, direct factorization gives
\begin{equation}
 H_{\mathrm{pl}}(t_*)^2-\frac45
 =\frac{(\omega-1)(5\omega-1)(25\omega^2+10\omega-11)}
 {20(2\omega-1)}.
\end{equation}
It is positive below $1/5$ and negative between $1/5$ and $7/15$.
On the shared planar branch,
\begin{equation}
 \{2P_3(c)\}^2-\frac45
 =\frac{(5\omega-3)(5\omega+3)^2}{40},
\end{equation}
which changes sign at $3/5$. The competing directions do not coincide
at either crossing. Formula \eqref{eq:transverse} supplies the middle branch.
\end{proof}

These intermediate directions have equal \emph{negative} overlaps
$e_1^\top u=e_2^\top u=-1/\sqrt5$, where $P_3$ is positive.
They exploit a secondary lobe of the overlap polynomial. The two reflected
directions cancel dependence that is odd under reflection across the teacher
plane when used with equal weights. This does not imply that an individual
student feature is confined to the teacher span.

\begin{theorem}[Legendre shared-branch exit]\label{thm:legendreexit}
For $\omega>3/5$, the single shared atom is globally optimal for
\begin{equation}\label{eq:legendreexit}
 \boxed{\frac\lambda{2P_3(c)}\leq\beta\leq
        \frac\lambda{5cs^2}.}
\end{equation}
It loses optimality through out-of-plane residual contacts
\begin{equation}
 u=-\tfrac12 e_+\pm\tfrac{\sqrt3}{2}e_0.
\end{equation}
The in-plane local splitting threshold is twice as large,
$\beta_\loc=2\lambda/(5cs^2)$, and is preempted.
\end{theorem}
\begin{proof}
Here $a=5/2,b=3/2$, $B=5\beta cs^2/\lambda$, and $B_c=2$.
Equation \eqref{eq:negative} satisfies $C_3(B)=1$ exactly at $B=1$:
the squared equality is $(B-1)(B+2)^2=0$. At $B=1$, the negative
minimum of \eqref{eq:residual} has $x=1/2,y=0$ and value $-1$;
reverse the direction for the positive contacts. For $B\leq1$ the
planar upper bound and transverse lower bound prove the global certificate.
For $B>1$ the transverse bound violates it. The interval is nonempty
exactly when $\omega\geq3/5$.
\end{proof}

\paragraph{Numerical scale at $\omega=0.8$.}
With $\lambda=1$, the Hermite interval is approximately
$[0.5856,2.6352]$, the $D=4$ Gegenbauer interval is
$[0.6588,3.2940]$, and the Legendre interval is
$[0.7027,2.1082]$. The Legendre in-plane local calculation alone
would instead predict an endpoint near $4.2164$.

\section{A many-teacher result: the modified spectral stability test}

The earlier constant-row-sum stability calculation also extends exactly.
Let $k\geq2$ unit teacher directions have Gram matrix $G$ with
$G\mathbf1=g\mathbf1$, $g>0$. Define
\begin{equation}
 u_0=\frac{\sum_i e_i}{\sqrt{kg}},\qquad
 c_0=e_i^\top u_0=\sqrt{\frac gk}.
\end{equation}
The sum direction is stationary for $H_D(u)=\sum_iK_D(e_i^\top u)$.
For a unit tangent vector $v\perp u_0$, differentiating along the great
circle $u(\theta)=u_0\cos\theta+v\sin\theta$ gives
\begin{equation}\label{eq:manyhessian}
 \left.\frac{d^2H_D(u(\theta))}{d\theta^2}\right|_0
 =K_D''(c_0)\sum_i(e_i^\top v)^2-kc_0K_D'(c_0).
\end{equation}
The nonzero eigenvalues of $\sum_i e_ie_i^\top$ on this tangent space
are those of $G$ on $\mathbf1^\perp$; ambient directions outside the
teacher span contribute eigenvalue zero. Since $K_D''(c_0)=6a_Dc_0>0$,
strict negative definiteness is equivalent to
\begin{equation}\label{eq:manystability}
 \boxed{2\lambda_{\max}\!\left(G\big|_{\mathbf1^\perp}\right)
       <g-\frac{k}{D+2}.}
\end{equation}
Including zero eigenvalues does not change the maximum because $G$ is
positive semidefinite. If the inequality is reversed, the sum direction
has a positive-curvature mode and is not a local maximum. At equality,
higher-order terms are needed. This recovers the Hermite criterion
$2\lambda_{\max}<g$ as $D\to\infty$ with fixed $k$.

For $k\leq D$ equicorrelated teachers with $0\leq\omega<1$,
$g=1+(k-1)\omega$ and the transverse Gram eigenvalue is $1-\omega$.
The strict local stability condition becomes
\begin{equation}\label{eq:manylocal}
 \boxed{\omega>\omega_\loc(k,D)
       =\frac{1+k/(D+2)}{k+1}.}
\end{equation}
For $k=2$, this agrees with the planar threshold in
\eqref{eq:tstar}. It is a global onset threshold for $D\geq4$ by the
earlier theorem, but not for $D=3$. For $k\geq3$, we have proved a
local stability boundary, not a global coexistence boundary. This is
the same distinction that mattered in the recent many-feature Hermite
calculation, where competing local maxima can exchange dominance before
either loses stability.

\section{The circle: orthogonal Fourier modes with no splitting}

For a second, nonpolynomial comparison, let $\varphi$ be uniform on the
circle and take the orthonormal Fourier family. A translated degree-$l$
feature is
\begin{equation}
 \phi_{l,\theta}(\varphi)=\sqrt2\cos(l(\varphi-\theta)),\qquad
 \ip{\phi_{l,\theta}}{\phi_{l,\vartheta}}=\cos(l(\theta-\vartheta)).
\end{equation}
Two equal teachers at angles $\pm\alpha$ satisfy the exact identity
\begin{equation}
 \beta\{\phi_{l,\alpha}+\phi_{l,-\alpha}\}
       =2\beta\cos(l\alpha)\phi_{l,0}.
\end{equation}
Their sum is already a single dictionary atom times a scalar. Its optimum
under \eqref{eq:measure} is the same atom with strength
$[2\beta|\cos(l\alpha)|-\lambda]_+$ and the appropriate sign.
Indeed $\norm{f_\mu}_{L^2}\leq\norm\mu_\TV$, and this representation
attains equality, so ordinary radial soft thresholding gives a global
minimizer. There is no need to split at any signal strength; if
$\cos(l\alpha)=0$, the target vanishes identically. This is also the
$D=2$ zonal-harmonic limit. It shows directly why orthogonality alone
does not enforce the Hermite phase diagram.

\section{What survives from the microscopic derivation}\label{sec:micro}

\subsection{The weight-prior geometry survives}
Changing the input law does not change an independently specified spherical
weight prior. For a fixed $m$-dimensional subspace $V$, the projection
$p=P_Vw$ of a uniform unit weight has density
\begin{equation}
 p_V(p)\propto(1-\norm p^2)^{(D-m-2)/2},\quad \norm p<1.
\end{equation}
Hence the leading cost per dimension for a fixed nonzero projection
length $r$ remains
\begin{equation}
 J(r)=-\tfrac12\log(1-r^2).
\end{equation}
This statement is independent of Hermite polynomials.

\subsection{The finite-dimensional committee output changes}
Take $D>m$, $w=ru+\sqrt{1-r^2}\,\eta$, with $u\in V$ unit and
$\eta$ uniform on the unit sphere of $V^\perp$. At fixed input $S$, set
$z=u^\top S$ and $Q=\norm{P_VS}^2$. Averaging over the transverse weights
gives, \emph{exactly},
\begin{align}\label{eq:committee}
 \E_\eta\phi_{3,w}(S)
 ={}&r^3\phi_{3,u}(S)\\
 &+\frac{3\sqrt{h_{3,D}}}{D-1}r(1-r^2)z
     \frac{m+2-(D+2)Q}{D-m}.\nonumber
\end{align}
To verify it, use $\E_\eta(\eta^\top S)^2=(1-Q)/(D-m)$ and vanishing
odd transverse moments in the cubic polynomial. In contrast to
\eqref{eq:hermite}, the second term usually does not vanish. Therefore
the exact finite-$D$ output of these committees is not a single
$r^3$-rescaled ridge atom, and one cannot simply import the old
$\lambda\propto\min_rJ(r)/r^3$ as an exact finite-dimensional result.
For $m=1$, the expression simplifies to
$K_D(r)\phi_{3,u}$, another check on the formula.

For fixed $m$ and fixed $r$, the correction in \eqref{eq:committee}
is $O(D^{-1})$ in $L^2$: $Q$ has beta distribution
$\operatorname{Beta}(m/2,(D-m)/2)$, $z=O_{L^p}(D^{-1/2})$
for fixed finite $p$, and the displayed factors give the bound.
With $X=\sqrt D S$, the normalized cubic satisfies
\begin{equation}
 \phi_{3,u}(S)=
 \sqrt{\frac{(D+2)^2(D+4)}{6D^2(D-1)}}
 \left\{(u^\top X)^3-\frac{3D}{D+2}u^\top X\right\}.
\end{equation}
It tends to the normalized Gaussian cubic. Thus the original radius
optimization and order-$D$ aligned-neuron scaling are consistent leading
candidates in the same high-dimensional regime, assuming the rest of the
reduction can be justified. This observation is not a proof of that reduction.

\subsection{What has not been transferred}
The paper integrates a weak background at width and sample size of order
$D^2$, and its effective data force may depend on residual variance.
The corresponding covariance law, determinant/free-energy terms, and any
self-consistency for $g$ must be derived again for spherical inputs.
Population risk is not automatically the full sample posterior in this
joint limit. The finite-$D$ Legendre theorems concern the explicitly defined
penalized network \eqref{eq:measure}, not order-$D$ committees in a proven
finite-sample posterior. Likewise, changing the readout prior, readout cap,
or total norm can change the population cost and its preferred sparsity.

\section{General activations, mixed bases, and a solvable mixture}

\subsection{The exact formulas for a general ridge student}
Under Gaussian inputs, expand a teacher activation $\tau$ and student
activation $\sigma$ in normalized Hermites:
\begin{equation}
 \tau(z)=\sum_{l\geq0}\alpha_l\psi_l(z),\qquad
 \sigma(z)=\sum_{l\geq0}\gamma_l\psi_l(z),
 \qquad \sum_l\alpha_l^2,\ \sum_l\gamma_l^2<\infty.
\end{equation}
Then the exact cross-kernel and student kernel are
\begin{equation}\label{eq:generalK}
 K_{TS}(q)=\sum_l\alpha_l\gamma_lq^l,
 \qquad K_{SS}(q)=\sum_l\gamma_l^2q^l.
\end{equation}
For $f_*=\beta\sum_i\tau(e_i^\top X)$, the linear term in
\eqref{eq:Fexpand} uses $\sum_iK_{TS}(e_i^\top u)$, whereas the quadratic
term uses $K_{SS}$. Orthogonality diagonalizes the bookkeeping; it does not
make the neuron weights for different degrees independent. For a
general multivariate teacher, simply use
$H(u)=\ip{f_*}{\sigma(u^\top X)}$ without a ridge decomposition.

The spherical version of \eqref{eq:generalK} replaces $q^l$ by
$Z_l^{(D)}(q)$ and uses expansion coefficients in the normalized
Gegenbauer basis. If teacher and student are written in different
polynomial bases, first express both in a common basis for the
\emph{actual input law}. The bases themselves do not determine a learning
law independently of the measure.

For the microscopic Gaussian committees, transverse averaging gives
\begin{equation}\label{eq:mixedradius}
 \E_Z\sigma(rz+\sqrt{1-r^2}Z)=\sum_l\gamma_l r^l\psi_l(z).
\end{equation}
Consequently the ratio of generated degrees changes with $r$. The original
single-radius optimization does not generally reduce to an unchanged
penalty on copies of $\sigma$; the atomic dictionary must include $r$, or
the cost must be derived jointly over its output coefficient vector.
Pairing signs cancels even degrees but does not separate odd degrees
from one another.

\subsection{An exact mixed linear--cubic phase diagram for onset and sharing}
For a concrete penalized-network calculation, take the same activation
for teacher and student,
\begin{equation}
 \sigma_\eta(z)=\sqrt\eta\,\psi_1(z)+\sqrt{1-\eta}\,\psi_3(z),
 \quad 0\leq\eta<1,
 \qquad K_\eta(q)=\eta q+(1-\eta)q^3.
\end{equation}
Here we mean actual unit-direction ridge neurons in
\eqref{eq:measure}, not the radius-optimized committees of
\eqref{eq:mixedradius}. Write $a=1-\eta$, $v=\eta$, and $\kappa=3-2\eta$.
The same cubic calculation gives
\begin{equation}\label{eq:mixedonset}
 \omega_c(\eta)=\frac13-\frac{2\eta}{9(1-\eta)}.
\end{equation}
For $\omega\geq0$, sharing is preferred at onset whenever
$\omega>\omega_c(\eta)$; if $\eta\geq3/5$, the shared direction is
optimal for every nonnegative correlation (with the marginal boundary
case at $\eta=3/5,\omega=0$). On the shared branch,
\begin{equation}\label{eq:mixedexit}
 \beta_\on=\frac\lambda{2K_\eta(c)},\qquad
 \boxed{\beta_\exit=
 \frac{\lambda(3-2\eta)}{4(1-\eta)(3+\eta)cs^2}.}
\end{equation}
These are global boundaries, not just local tests. To see this, put
$T=2\beta a cs^2/\lambda$. The shared residual is
\begin{equation}
 \frac{\ip R{\sigma_\eta(u^\top X)}}\lambda
   =(a-vT)x^3+3Txy^2+v(1+T)x.
\end{equation}
For $x\geq0$, it is nonnegative at $y=0$, and maximized on the teacher
plane. There it equals
$[a-(v+3)T]x^3+[v+(v+3)T]x$, is one at $x=1$, and is bounded above
by one precisely until $T=\kappa/[2(v+3)]$. The residual certificate
therefore proves \eqref{eq:mixedexit}. Differentiating the target overlap
proves \eqref{eq:mixedonset}. In the pure-linear limit $\eta=1$, the
target is a single linear ridge and no split is needed.

\paragraph{Putting a Legendre polynomial on Gaussian inputs.}
The unnormalized polynomial satisfies
\begin{equation}
 P_3(z)=\frac{5\sqrt6}{2}\psi_3(z)+6\psi_1(z).
\end{equation}
Its variance-normalized Gaussian kernel is
$K(q)=(25/49)q^3+(24/49)q$. Thus the corresponding constant-penalty,
matched teacher/student onset threshold is $\omega_c=3/25$, not
$1/3$, and not the spherical Legendre value $3/5$. This is a degree
mixture under Gaussian inputs, not an alternative Gaussian orthogonal family.

\subsection{Local tests for any smooth normalized matched kernel}
Let $K(1)=1$, use the same feature family for teachers and students, and
assume the derivatives below exist. At the shared direction,
\begin{equation}
 m=2\beta K(c)-\lambda,
 \qquad \text{local onset stability: } cK'(c)-s^2K''(c)\geq0.
\end{equation}
The reflected-pair ansatz has cost
\begin{equation}
 F(m,t)=\lambda m-\beta m H_K(t)
 +\frac{m^2}{4}\{1+K(2t^2-1)\},
\end{equation}
where $H_K(t)=K(ct+s\sqrt{1-t^2})+K(ct-s\sqrt{1-t^2})$.
When the denominator is positive, its local shared-splitting threshold is
\begin{equation}\label{eq:generalsplit}
 \beta_\loc=
 \frac{\lambda K'(1)}
 {2\{K(c)K'(1)-cK'(c)+s^2K''(c)\}}.
\end{equation}
To derive this, expand in the separation angle $\theta$ with
$t=\cos\theta$. The coefficient of $\theta^2$ is
$\beta m[cK'(c)-s^2K''(c)]-m^2K'(1)/2$.
The Legendre example proves why \eqref{eq:generalsplit} cannot generally
be promoted to a global phase boundary without the residual test.

\section{More realistic students: a concrete obstruction and a usable route}

\subsection{Bias-free ReLU cannot learn these pure cubic teachers}
For symmetric Gaussian inputs,
\begin{equation}
 \operatorname{ReLU}(z)=\tfrac12z+\tfrac12|z|.
\end{equation}
Its odd component is purely linear. Every finite or square-integrable
infinite linear combination of bias-free ReLU ridges therefore has an
odd component of degree one; its remaining part is even. A pure cubic
Hermite teacher is orthogonal to both parts. Consequently,
\begin{equation}
 \ip{f_*}{f_{\mathrm{ReLU}}}=0,
 \qquad \norm{f_*-f_{\mathrm{ReLU}}}^2
   =\norm{f_*}^2+\norm{f_{\mathrm{ReLU}}}^2.
\end{equation}
The population optimum is the zero function, even without a penalty.
Thus merely replacing the cubic student by a familiar activation need not
preserve any cubic-learning phase. This obstruction applies to a single
hidden layer with no hidden biases under the symmetric input measure;
it is not a claim about general ReLU networks.

Hidden biases remove this particular obstruction. For
$\sigma_h(z)=\operatorname{ReLU}(z+h)$, its third normalized Hermite
coefficient is
\begin{equation}
 \gamma_3(h)=\E[\sigma_h(Z)\psi_3(Z)]
       =-\frac{h\,\varphi(h)}{\sqrt6},
\end{equation}
where $\varphi$ is the standard Gaussian density. This follows by Gaussian
integration by parts, or direct one-dimensional integration. For $h\ne0$
the cubic channel is present, together with other degrees. The feasible
neuron becomes an atom indexed by $(u,h)$, and its normalization and penalty
can depend on $h$. Treating only $\gamma_3$ while dropping the other
generated degrees would miss the fit penalty and interactions.

For matched ReLU teachers and students, an exact Gaussian overlap kernel
is available as well; for $\sigma(z)=\sqrt2\operatorname{ReLU}(z)$,
\begin{equation}
 K(q)=\frac{\sqrt{1-q^2}+(\pi-\arccos q)q}{\pi}.
\end{equation}
This is the arc-cosine kernel \cite{cho}. Because ReLU is not odd, signed
readouts must be retained explicitly, and the nonzero mean channel must
also be kept or deliberately centered. The general certificate still
applies; a cubic-only calculation does not.

\subsection{Non-Gaussian and manifold inputs}
For an arbitrary input law $P_X$, the exact objects are
\begin{equation}
 K(u,v)=\E_{P_X}[\sigma(u^\top X)\sigma(v^\top X)],\qquad
 h(u)=\E_{P_X}[f_*(X)\sigma(u^\top X)].
\end{equation}
They may depend on the full directions, not only $u^\top v$.
Equations \eqref{eq:measure} and \eqref{eq:certificate} remain true, but the
one-variable reductions need not. With latent manifold inputs $X=g(Z)$,
one can expand the actual composed neuron $\sigma(u^\top g(Z))$ in an
orthonormal basis for $Z$ and form these inner products from its coefficients.
This is exact bookkeeping if the full expansion is retained, not an
automatic proof of the original Gaussian posterior theory.

\section{What generalizes, and how strongly}

\paragraph{Exact structure.}
For an explicitly specified linear population penalty and squared loss,
the residual inequality, onset test, interaction Gram matrix, and
optimization over measures are independent of Hermite polynomials.
An ordinary two-layer network with optimized directions is already in this
framework. The real restriction is the form of its penalty and the tractable
overlap kernel, not the use of the words teacher and student.

\paragraph{Robustness under small changes.}
If two normalized matched kernels obey
$\sup_{q\in[-1,1]}|\widetilde K(q)-K(q)|\leq\varepsilon$, then for
$k$ equal teachers
\begin{equation}
 \sup_u|\widetilde H(u)-H(u)|\leq k\varepsilon.
\end{equation}
If a maximizer has a height gap $\delta$ over directions outside a chosen
neighborhood, that localization is preserved whenever
$2k\varepsilon<\delta$. Here
\begin{equation}
 \sup_{q\in[-1,1]}|K_D(q)-q^3|
      =\frac{2}{\sqrt3(D-1)}.
\end{equation}
This gives a concrete high-dimensional robustness statement for any fixed
number of teachers. A nondegenerate local maximum persists under small
$C^2$ perturbations; two such branches crossing transversely give a
persistent, shifted first-order transition. Degenerate boundaries and
vanishing gaps require a separate analysis. In particular, the earlier
all-$k$ Hermite equicorrelation formula is not an exact Gegenbauer formula
just because the kernels are close.

\paragraph{What is model-dependent.}
The exact $1/3$ threshold, the assumption that all optimal atoms lie in the
teacher span, the endpoint of the shared branch, and even the existence of
a splitting phase depend on the kernel and feasible dictionary. The
Legendre and Fourier calculations give explicit counterexamples to a claim
of universality based only on orthogonality. Neuron counts, readout
saturation, and the distributed-versus-aligned distinction also depend on
the microscopic prior and asymptotic scaling.

\paragraph{What remains open in this extension.}
We have not classified the post-exit populations for the spherical model,
proved its finite-sample Bayesian background integration, or established
that SGD reaches the population optimum. We also have not redone the
all-$k$ equicorrelation theorem for Gegenbauer atoms. The present results
fully resolve two-teacher onset and the shared-state interval, and isolate
specific places where a broader theory needs new work.

\paragraph{A practical next model.}
A useful continuation is a two-layer biased-ReLU student on Gaussian
inputs, with the existing cubic teacher and an explicitly specified weight
penalty. Evaluate the cross-kernel and student kernel using the full
Hermite spectrum or accurate quadrature; optimize over direction and bias;
then certify the resulting populations with the residual inequality. A
separate finite-width experiment can compare that optimum with trained
networks. Only after that comparison should one add a non-Gaussian or
latent-manifold input law, so changes from the activation, optimization,
and input geometry can be identified separately.

\section{Verification and provenance}
The accompanying reproducible calculation, \texttt{verify.py}, checks
the cubic spherical covariance by exact symbolic moments and checks the
Legendre transition factorizations. It also integrates the Legendre
covariance by deterministic spherical quadrature and compares the analytic
onset values with multi-start optimization over the full direction sphere.
The residual is checked inside, at, and beyond the predicted shared
endpoint for several dimensions and correlations. Numerical optimization
is a corroborating check; the global claims follow from the proofs above.
Across 60 choices of dimension and correlation, analytic onset maxima
agree with the full-sphere searches to below $10^{-12}$; deterministic
Legendre covariance quadrature agrees to below $10^{-12}$ as well. The
Gaussian degree-mixture formulas and the biased-ReLU cubic coefficient
are checked separately. The source material was read without modification.

\begin{thebibliography}{9}
\bibitem{sources}
Project reference documents, \emph{Statistical Physics of Internal
Representations in Feature Learning}, September 17, 2026, especially
Appendix E; and \emph{Correlated features and shared neurons}, sections 1--6.
Also consulted the related chats on correlated features, generalized
spherical harmonics, and manifold inputs.
\bibitem{dlmf}
NIST Digital Library of Mathematical Functions, Section 18.38(ii),
\href{https://dlmf.nist.gov/18.38}{Zonal Spherical Harmonics}.
\bibitem{frye}
C. Frye and C. Efthimiou, \emph{Spherical Harmonics in $p$ Dimensions},
2012. \href{https://arxiv.org/abs/1205.3548}{arXiv:1205.3548}.
\bibitem{bach}
F. Bach, \emph{Breaking the Curse of Dimensionality with Convex Neural
Networks}, JMLR 18(19), 1--53, 2017.
\href{https://jmlr.org/papers/v18/14-546.html}{Publisher's article}.
\bibitem{cho}
Y. Cho and L. K. Saul, \emph{Kernel Methods for Deep Learning},
NeurIPS 22, 2009.
\href{https://papers.nips.cc/paper/2009/file/5751ec3e9a4feab575962e78e006250d-Paper.pdf}{Original paper}.
\end{thebibliography}
\end{document}
